\documentclass{article}

\usepackage{booktabs}
\usepackage{tabularx}

\title{Development Plan\\\progname}

\author{\authname}

\date{}

\input{../Comments.text}
\input{../Common.text}

\begin{document}

\maketitle

\begin{table}[hp]
\caption{Revision History} \label{TblRevisionHistory}
\begin{tabularx}{\textwidth}{llX}
\toprule
\textbf{Date} & \textbf{Developer(s)} & \textbf{Change}\\
\midrule
Sept 22, 2026 & Team 11 & Initial draft for TA meeting\\
\bottomrule
\end{tabularx}
\end{table}

\newpage{}

This document lays out how Team 11 plans to work together over the course of
the Tandem capstone project: our meeting and communication plans, team
roles, git/GitHub workflow, project scheduling, our Proof of Concept
demonstration plan, expected technology stack, and our team charter.

\section{Confidential Information?}

This project does not involve any confidential information from industry.
The team does not currently have a direct industry supervisor, so there is no
NDA or confidentiality agreement in place.

\section{IP to Protect}

There is no IP to protect at this stage. This is an open-source systems
research project (extending existing open-source LLM inference engines such
as vLLM/SGLang) with no industry sponsor, so no IP agreement or
Intellectual Property Guide Acknowledgement is required.

\section{Copyright License}

The team is adopting the MIT License, included as \texttt{LICENSE} in the
root of the
\href{https://github.com/jjustin-k/Tandem}{project repository}.

\section{Team Meeting Plan}

\todo{CONFIRM WITH TEAM: exact day/time and virtual vs. in-person.}

The team will meet weekly (proposed: once during a shared lecture/tutorial
slot when nothing else is scheduled, day/time TBD) to review progress, and
more often as needed leading up to deliverable deadlines. Meetings will
primarily be in person on campus, with virtual meetings (Discord/Zoom, TBD)
used when schedules conflict.

Each meeting will have a rotating chair responsible for setting the agenda
ahead of time and keeping the meeting on track, and a notetaker who records
decisions and action items as a GitHub Issue or in the meeting minutes.

We do not yet have a supervisor/advisor confirmed (actively in discussion
with a few candidates, per our project review feedback). Once one is
secured, we plan to meet with them roughly biweekly, virtually by default,
with cadence adjusted to their availability.

\section{Team Communication Plan}

\todo{CONFIRM WITH TEAM: which chat tool (Discord/Slack/Teams).}

Day-to-day communication will happen over a team chat tool (Discord,
tentative). GitHub Issues will be used for anything tied to concrete work:
bug reports, feature/research tasks, and decisions that affect the code or
documentation, so that discussion stays attached to the artifact it affects
and is searchable later.

\section{Team Member Roles}

At this stage we anticipate the following rotating roles rather than fixed
assignments, since responsibilities will naturally shift as the project
moves from research/prototyping into systems integration and benchmarking:

\begin{itemize}
  \item \textbf{Meeting chair} --- sets the agenda, runs the meeting.
  \item \textbf{Notetaker} --- records decisions/action items.
  \item \textbf{PR reviewer(s)} --- at least one teammate other than the
  author reviews each pull request before merge.
  \item \textbf{Liaison} --- Justin Kwinecki is the team liaison for the
  course (point of contact with the instructor/TA).
\end{itemize}

\section{Workflow Plan}

\begin{itemize}
	\item Git/GitHub: work happens on feature branches off \texttt{main}
	(named \texttt{<initials>/<short-description>}); changes land via pull
	request with at least one team member's approval before merging. No
	direct pushes to \texttt{main} for code changes.
	\item Issues: GitHub Issues are used for tracking bugs, features, and
	research tasks, labelled by type (\texttt{bug}, \texttt{feature},
	\texttt{research}, \texttt{docs}) and by deliverable/milestone.
	\item CI/CD: GitHub Actions runs on every pull request to build the
	LaTeX documentation (already set up, see
	\texttt{.github/workflows/latex-pages.yml}) and, once source code
	exists, will run linting and the test suite.
\end{itemize}

\section{Project Decomposition and Scheduling}

\todo{CONFIRM WITH TEAM / ADD LINK: GitHub Projects board URL.}

\begin{itemize}
  \item We will use a GitHub Projects board (link TBD) with columns for
  Backlog, In Progress, In Review, and Done, with issues tied to each
  deliverable milestone below.
  \item Big-picture schedule follows the course deliverable calendar
  (Problem Statement \& Dev Plan, Proof of Concept demo, Requirements/SRS,
  Hazard Analysis, V\&V Plan, Design Documents, Rev 0 demo, Final demo, and
  V\&V Report), interleaved with our own technical milestones:
  \begin{itemize}
    \item \textbf{Months 1--3}: Prove the fused prompt-tree prefill
    mechanism in pure PyTorch (no serving engine) --- a shared trunk plus
    multiple branch prompts evaluated in a single forward pass with a
    custom tree-attention mask, producing identical logits to running each
    branch separately.
    \item \textbf{Months 4--5}: Prototype the ``KV-Promise'' scheduler
    extension (binding staggered concurrent requests to an in-progress
    prefill) on top of a real serving engine (vLLM).
    \item \textbf{Months 6--8}: End-to-end integration with a multi-agent
    framework and benchmarking (Time-To-First-Token, VRAM, FLOPs) against
    baseline vLLM across fan-out factors.
  \end{itemize}
\end{itemize}

\section{Proof of Concept Demonstration Plan}

The main risk to this project is technical: whether a single fused
attention pass can correctly and efficiently evaluate a shared-prefix
``prompt tree'' (one trunk plus several divergent branches) instead of
recomputing the shared prefix once per branch, and whether that mechanism
still holds up once integrated into a real serving engine under concurrent
load.

For the Proof of Concept demonstration, we will show a pure PyTorch
implementation that takes a shared trunk plus several branch prompts,
packs them into one sequence, builds a custom 2D block-structured
tree-attention mask, and runs a single forward pass. We will demonstrate
that the resulting logits for each branch are numerically identical to
running each branch as its own independent forward pass --- proving the
core mechanism is correct before we invest in performance engineering or
vLLM integration.

\section{Expected Technology}

\begin{itemize}
\item \textbf{Language}: Python, with C++/CUDA if custom kernel work is
needed later.
\item \textbf{Libraries/frameworks}: PyTorch, Hugging Face Transformers,
FlashAttention, and (later) a fork of vLLM's scheduler/\texttt{BlockManager}.
\item \textbf{Pre-trained models}: we will use existing open-weight LLMs
(e.g. Llama/Qwen family) rather than training our own model --- this
project is about serving-system performance, not model quality.
\item \textbf{Linting}: \texttt{ruff}/\texttt{black} for Python.
\item \textbf{Testing}: \texttt{pytest}, comparing fused-pass logits
against a separate-pass baseline for correctness.
\item \textbf{Code coverage}: \texttt{coverage.py} (to be investigated
further once the codebase exists).
\item \textbf{CI}: GitHub Actions (already building/deploying
documentation; will be extended to run lint/tests).
\item \textbf{Performance measurement}: NVIDIA Nsight Systems and the
PyTorch Profiler, benchmarking on the department's GPU clusters
(\texttt{ada}: 12$\times$ L40s for development, \texttt{grace}: 4$\times$
H100 for benchmarking).
\item Git, GitHub, and GitHub Projects for version control and project
tracking.
\end{itemize}

\section{Use of AI and LLMs}

\todo{CONFIRM WITH TEAM before submitting.}

The team is using LLM coding assistants (Claude Code, and potentially
GitHub Copilot) to accelerate boilerplate code, draft and structure
documentation, and research related systems literature (e.g. comparing our
approach against vLLM/SGLang/Medusa). AI-generated code and text is
reviewed by a team member before being merged; the team verifies technical
claims (especially about the internals of existing systems like vLLM)
against primary sources rather than trusting LLM output directly. AI is
not used to make architectural or grading-relevant decisions without team
review.

\section{Coding Standard}

Python code will follow PEP 8, enforced via \texttt{ruff}/\texttt{black},
with type hints and docstrings on public functions. This will be
revisited once we know whether C++/CUDA code is needed.

\newpage{}

\section{Appendix --- Generative AI Use Disclosure}

\subsection{AI Tools Used}

Claude (Claude Sonnet 5), used via the Claude Code CLI.

\subsection{How and Where Used}

Claude was used to draft the initial structure and text of this Development
Plan (Sections 1--10 and this appendix) from information and decisions
provided by the team, and to help brainstorm and refine the technical
approach of the Tandem project itself, including identifying which parts of
the design are genuinely novel versus already implemented by existing
systems such as vLLM/SGLang.

\subsection{Verification of AI-Generated Content}

\todo{CONFIRM WITH TEAM before submitting: describe the team's actual
review pass.}

The team reviewed all AI-drafted text before including it in this document
and will revise sections (meeting logistics, roles, schedule, etc.) to
reflect the team's actual decisions rather than AI-proposed defaults.
Technical claims about existing systems (e.g. vLLM's prefix caching and
Copy-on-Write behaviour) were checked against those projects' own
documentation. The team is responsible for the accuracy, completeness,
consistency, and appropriateness of this document, regardless of whether
material was AI- or human-generated.

\newpage{}

\section*{Appendix --- Reflection}

\wss{Not required for CAS 741}

\input{../Reflection.text}

\todo{CONFIRM WITH TEAM: draft answers below, replace with your own.}

\begin{enumerate}
    \item A development plan forces the team to make explicit decisions
    (workflow, communication, roles, schedule) before deadline pressure
    hits, so disagreements about ``how we work'' get resolved once instead
    of repeatedly derailing later deliverables.
    \item CI/CD catches build/documentation errors and (later) test
    failures before they reach \texttt{main}, and automates otherwise
    tedious tasks like PDF generation and deployment. The main disadvantage
    is upfront setup/maintenance cost, and that a flaky or slow pipeline
    can become a bottleneck if not kept reliable.
    \item Draft: no major disagreements yet at this stage; the team was
    initially considering a broader, more general project before narrowing
    scope to the two specific systems gaps described in this plan, which
    was a scope discussion rather than a disagreement.
\end{enumerate}

\newpage{}

\section*{Appendix --- Team Charter}

\wss{borrows from
\href{https://engineering.up.edu/industry_partnerships/files/team-charter.pdf}
{University of Portland Team Charter}}

\subsection*{External Goals}

\todo{CONFIRM WITH TEAM.}

Draft: build something substantial enough to be a genuine portfolio/interview
talking point in systems engineering, and aim for a strong final grade;
stretch goal of presenting at the Capstone EXPO if results are compelling.

\subsection*{Attendance}

\subsubsection*{Expectations}

Team members are expected to attend all scheduled team meetings on time and
stay for the full meeting; if unable to attend, notify the team as far in
advance as possible.

\subsubsection*{Acceptable Excuse}

Illness, other course/exam conflicts, and emergencies are acceptable
excuses. Simply forgetting, or deprioritizing without advance notice, is
not.

\subsubsection*{In Case of Emergency}

Notify the team chat and the meeting chair as soon as possible; the
notetaker will bring the absent member up to speed on decisions/action
items afterward.

\subsection*{Accountability and Teamwork}

\subsubsection*{Quality}

Members are expected to come to meetings having completed their agreed-upon
work, or to flag blockers before the meeting rather than at it. Pull
requests should be self-reviewed (tests passing, docs updated) before
requesting review.

\subsubsection*{Attitude}

We will adopt the
\href{https://www.contributor-covenant.org/}{Contributor Covenant} as our
code of conduct. Disagreements are expected to be resolved through
discussion and, if unresolved, a team vote (see Decision Making below).

\subsubsection*{Stay on Track}

\todo{CONFIRM WITH TEAM: target metrics and consequences.}

We will use the course's project management metrics (commits, issue
activity) as a lightweight check-in signal. If a member consistently
under-contributes, the team will raise it directly with them first, and
escalate to the TA/instructor only if it persists.

\subsubsection*{Team Building}

\todo{CONFIRM WITH TEAM.}

\subsubsection*{Decision Making}

Default to consensus; if consensus isn't reached after discussion, decide
by majority vote among team members.

\end{document}